\documentclass[11pt]{article}
\usepackage[margin=0.95in]{geometry}
\usepackage{amsmath,amssymb,amsthm}
\usepackage[colorlinks=true,linkcolor=blue,citecolor=blue,urlcolor=blue]{hyperref}
\newtheorem{theorem}{Theorem}[section]
\newtheorem{lemma}[theorem]{Lemma}
\newtheorem{proposition}[theorem]{Proposition}
\newtheorem{corollary}[theorem]{Corollary}
\theoremstyle{remark}\newtheorem{remark}[theorem]{Remark}
\newcommand{\m}{\mathfrak m}
\newcommand{\n}{\mathfrak n}
\newcommand{\gr}{\operatorname{gr}}
\newcommand{\edim}{\operatorname{edim}}
\newcommand{\codim}{\operatorname{codim}}
\newcommand{\dx}{\operatorname{dx}}
\newcommand{\Soc}{\operatorname{Soc}}
\newcommand{\Dsg}{\mathsf D_{\mathrm{sg}}}
\newcommand{\thick}{\operatorname{thick}}
\title{Dominance and common deformations of short local rings}
\author{Henry Zweiman}
\date{September 9, 2026}
\begin{document}
\maketitle
\begin{abstract}
Let $(R,\m,k)$ be a commutative Artinian local ring with $\m^3=0$. We construct a one-dimensional complete Cohen--Macaulay local ring whose quotients by two regular elements outside the square of its maximal ideal are $R$ and $\gr_{\m}R$. Consequently, dominance and uniform dominance are each equivalent for $R$ and its associated graded ring. Their dominant indices differ by at most the bounds $\dx(R)\le4\dx(\gr_{\m}R)+3$ and the reverse inequality when finite. Combining this construction with Kimura's equicharacteristic dominance theorem and known deformation and flat-descent results, we prove that every non-complete-intersection Gorenstein local ring satisfying either $\m^3=0$ or $e(R)\le\codim R+2$ is dominant. This answers both clauses of Kobayashi--Takahashi Question 7.1 without restrictions on characteristic or residue-field cardinality.
\end{abstract}

\section{Introduction}
All rings in this paper are commutative, Noetherian, and have identity. A local ring $(R,\m,k)$ is \emph{dominant} if
\[
 k\in\thick_{\Dsg(R)}(X)\qquad\text{for every nonzero }X\in\Dsg(R).
\]
Here $\Dsg(R)$ is the quotient of the bounded derived category of finitely generated modules by the perfect complexes. A thick subcategory is closed under shifts, cones, and direct summands. Dominance was introduced by Takahashi in connection with subcategory classification \cite{T}. For an Artinian ring, $k$ generates $\Dsg(R)$, so dominance says that this category has no nonzero proper thick subcategory.

Short local rings, here meaning those with $\m^3=0$, have long provided examples and tests for homological conjectures; see Avramov--Iyengar--\c Sega \cite{AIS}. In the Gorenstein case, Kobayashi and Takahashi ask whether every non-complete-intersection ring with $\m^3=0$, or more generally with $e(R)\le\codim R+2$, is dominant \cite[Question 7.1]{KT}. The multiplicity is Hilbert--Samuel multiplicity, and $\codim R=\edim R-\dim R$.

Kimura recently proved the short case when $R$ contains a field, and the higher-dimensional case in equal characteristic with infinite residue field \cite[Theorem 7.3]{K}. His theorem supplies the essential representation-theoretic input here. In particular, the first non-complete-intersection Gorenstein dominant examples already belong to that work. Our contribution is a common deformation that transfers dominance from the associated graded ring to an arbitrary short local ring, including rings without a coefficient field. The construction does not require the Gorenstein hypothesis.

\begin{theorem}\label{thm:deform}
Let $(R,\m,k)$ be an Artinian local ring with $\m^3=0$, and put $G=\gr_{\m}R$. There exist a one-dimensional complete Cohen--Macaulay local ring $(T,\n,k)$ and $T$-regular elements $u,v\in\n\setminus\n^2$ such that
\begin{equation}\label{eq:quotients}
 T/(u)\cong G,\qquad T/(v)\cong R.
\end{equation}
If $R$ has residue characteristic $p>0$ and $p\ne0$ in $R$, one can take $T$ finite free of rank $\ell_R(R)$ over a complete Cohen discrete valuation ring $V$ with uniformizer $p$, with $u=p$ and $v=p-a$ for an element $a$ satisfying $a^3=0$.
\end{theorem}

The elements $u$ and $v$ need not be distinct. The condition that both lie outside $\n^2$ is essential for the categorical application: dominance need not descend through a regular element in $\n^2$; see \cite[Corollary 6.12]{T}.

\begin{theorem}\label{thm:transfer}
With $R$ and $G$ as above, $R$ is dominant if and only if $G$ is dominant. Also, $R$ is uniformly dominant if and only if $G$ is uniformly dominant. If their dominant indices are finite, then
\begin{equation}\label{eq:indices}
 \dx(R)\le4\dx(G)+3,\qquad \dx(G)\le4\dx(R)+3.
\end{equation}
\end{theorem}

\begin{theorem}\label{thm:gorenstein}
Let $(R,\m,k)$ be a Gorenstein local ring that is not a complete intersection. If either $\m^3=0$ or $e(R)\le\codim R+2$, then $R$ is dominant.
\end{theorem}

Thus Theorem~\ref{thm:gorenstein} answers Question 7.1 of \cite{KT}. The removal of the characteristic restriction uses Theorem~\ref{thm:deform}; the removal of the finite-residue-field obstruction in positive dimension uses Liu's flat descent \cite{L}. We make no assertion that the rings in Theorem~\ref{thm:gorenstein} are uniformly dominant. Theorem~\ref{thm:transfer} transfers that stronger question to their graded Artinian counterparts without answering it.

\section{The common deformation}
Write
\[
 e=\dim_k(\m/\m^2),\qquad f=\dim_k\m^2,\qquad N=1+e+f=\ell_R(R).
\]
Since $\m^3=0$, multiplication gives a well-defined symmetric bilinear map
\begin{equation}\label{eq:tensor}
 b:(\m/\m^2)\times(\m/\m^2)\longrightarrow\m^2.
\end{equation}
Its values span $\m^2$. The graded ring $G$ is precisely the vector space $k\oplus k^e\oplus k^f$ with multiplication given by $b$ in positive degree and all products of total degree at least three zero.

\begin{lemma}\label{lem:field}
If $R$ contains a field, then $R\cong G$ as local rings.
\end{lemma}
\begin{proof}
The Artinian ring $R$ is complete, so the equicharacteristic Cohen structure theorem gives a coefficient field identified with $k$; see \cite[Tag 032A]{Stacks}. Choose a $k$-linear complement $E$ of $\m^2$ in $\m$. The decomposition $R=k\oplus E\oplus\m^2$ has exactly the multiplication described in \eqref{eq:tensor}, because $E\m^2=(\m^2)^2=0$. Identifying $E$ with $\m/\m^2$ gives the desired isomorphism.
\end{proof}

We now treat rings that do not contain a field. Their residue field has characteristic $p>0$, with $p\ne0$ in $R$. By the Cohen structure theorem there is a complete discrete valuation ring $V$ of characteristic zero, with maximal ideal $pV$ and residue field $k$, and a local homomorphism
\begin{equation}\label{eq:cohen}
 V\longrightarrow R
\end{equation}
inducing the chosen residue-field isomorphism \cite[Tag 032A and its proof]{Stacks}. This statement holds for imperfect $k$ as well. The map \eqref{eq:cohen} is not asserted to be injective. All uses of scalars from $V$ in $R$ are through this map.

Choose elements $x_1,\ldots,x_e\in\m$ whose classes are a basis of $\m/\m^2$, and a $k$-basis $z_1,\ldots,z_f$ of $\m^2$. Define $b_{ij}^t\in k$ by
\[
 x_ix_j=\sum_{t=1}^f b_{ij}^t z_t.
\]
This equality uses the canonical $k$-module structure on $\m^2$, annihilated by $\m$. Choose symmetric lifts $B_{ij}^t=B_{ji}^t\in V$ of these coefficients. Empty lists and empty sums are allowed when $f=0$.

Define a free $V$-module
\begin{equation}\label{eq:algebra}
 T=V\cdot1\ \oplus\ \bigoplus_{i=1}^e VX_i\ \oplus\ \bigoplus_{t=1}^f VZ_t
\end{equation}
with a $V$-bilinear multiplication, unit $1$, and rules
\begin{equation}\label{eq:rules}
 X_iX_j=\sum_{t=1}^f B_{ij}^t Z_t,\qquad X_iZ_t=0,\qquad Z_sZ_t=0.
\end{equation}
These rules define an associative commutative algebra: on triples of nonunit basis elements, both associated products are zero; triples containing $1$ satisfy associativity automatically. Put
\[
 J=\bigoplus_i VX_i\oplus\bigoplus_t VZ_t.
\]
Then $J$ is an ideal with $J^3=0$ and $T/J=V$. Therefore $T$ is local, with maximal ideal $\n=pT+J$, and has dimension one. It is Noetherian and complete: it is finite free over $V$, and the $p$-adic and $\n$-adic topologies agree since $p^rT\subseteq\n^r$ and $\n^{r+2}\subseteq p^rT$. Multiplication by $p$ is injective, so $T$ is Cohen--Macaulay. Reducing \eqref{eq:rules} modulo $p$ gives
\begin{equation}\label{eq:fiber}
 T/pT\cong G.
\end{equation}

The coefficient map and the assignments $X_i\mapsto x_i$, $Z_t\mapsto z_t$ define a ring homomorphism $\varphi:T\to R$. Indeed, the differences between the lifted coefficients and their residues annihilate $\m^2$, and all products involving $\m^2\m$ vanish. It is surjective: first represent an element of $R$ modulo $\m$ by an element of $V$, then represent the remaining class modulo $\m^2$ by a $V$-linear combination of the $x_i$, and finally represent the remainder by the $z_t$.

Adapt the bases to the nonzero element $p\in R$. If $p\notin\m^2$, choose $x_1=p$ and set $a=X_1$. If $p\in\m^2$, choose $z_1=p$ and set $a=Z_1$. In both cases $a^3=0$ and $\varphi(p-a)=0$.

\begin{lemma}\label{lem:length}
The element $p-a$ is $T$-regular, the map $\varphi$ induces an isomorphism $T/(p-a)\cong R$, and both $p$ and $p-a$ belong to $\n\setminus\n^2$.
\end{lemma}
\begin{proof}
Let $K$ be the fraction field of $V$. In $T\otimes_V K$, the element $p-a$ is invertible, with inverse
\[
 p^{-1}\bigl(1+p^{-1}a+p^{-2}a^2\bigr).
\]
Since $T$ is $V$-free, it embeds in this algebra, proving that $p-a$ is regular.

In the ordered basis $(1,X_1,\ldots,X_e,Z_1,\ldots,Z_f)$, multiplication by $a$ is strictly triangular. Thus multiplication by $p-a$ has determinant $p^N$. The length formula for the cokernel of a square matrix of nonzero determinant over a discrete valuation ring gives
\[
 \ell_V\bigl(T/(p-a)\bigr)=v_p(p^N)=N.
\]
The $R$-composition factors of $R$ are all $k$, which also has $V$-length one. Hence $\ell_V(R)=N$. The surjection $T/(p-a)\to R$ is an isomorphism by additivity of length.

Finally, the projection $\pi:T\to V$ with kernel $J$ satisfies $\pi(\n^2)=p^2V$. Both $p$ and $p-a$ project to $p\notin p^2V$, so neither lies in $\n^2$.
\end{proof}

\begin{proof}[Proof of Theorem~\ref{thm:deform}]
When $R$ does not contain a field, the algebra \eqref{eq:algebra} and Lemma~\ref{lem:length} give the result with $u=p$ and $v=p-a$. In the remaining case, Lemma~\ref{lem:field} allows $T=R[[t]]$ and $u=v=t$. This ring is complete of dimension one, $t$ is regular and outside its maximal-ideal square, and an Artinian ring is Cohen--Macaulay.
\end{proof}

\section{Transfer of categorical generation}
We recall the exact deformation properties needed. If $(A,\mathfrak a,k)$ is local and $x\in\mathfrak a$ is $A$-regular, then
\begin{equation}\label{eq:known}
 A/(x)\text{ dominant}\ \Longrightarrow\ A\text{ dominant}.
\end{equation}
The converse holds when $x\notin\mathfrak a^2$ \cite[Theorem 5.6]{T}.

For quantitative generation, let $\langle X\rangle_j$ be the subcategory obtained from shifts, finite direct sums, and direct summands of $X$ using at most $j-1$ successive extensions, with $\langle X\rangle_0=\{0\}$. Takahashi's dominant index is
\[
 \dx(A)=\inf\{r\in\mathbb Z_{\ge-1}:k\in\langle X\rangle_{r+1}
       \text{ for every }0\ne X\in\Dsg(A)\}.
\]
A ring is \emph{uniformly dominant} when this index is finite. A regular ring has index $-1$. For a regular element $x$, the known bounds are
\begin{equation}\label{eq:knownindices}
 \dx(A)\le2\dx(A/(x))+1,
 \qquad
 \dx(A/(x))\le2\dx(A)+1\quad\text{if }x\notin\mathfrak a^2;
\end{equation}
see \cite[Theorem 6.2]{U}.

\begin{proof}[Proof of Theorem~\ref{thm:transfer}]
Choose $T,u,v$ from Theorem~\ref{thm:deform}. Apply \eqref{eq:known} and its converse to each of $u,v$. This gives
\[
 G\text{ dominant}\ \Longleftrightarrow\ T\text{ dominant}
 \ \Longleftrightarrow\ R\text{ dominant}.
\]
Applying \eqref{eq:knownindices} twice gives
\[
 \dx(R)\le2\dx(T)+1\le4\dx(G)+3.
\]
Interchanging $u$ and $v$ gives the other inequality and the equivalence of finiteness. No uniform bound is inferred from dominance alone.
\end{proof}

\section{Gorenstein rings of almost minimal multiplicity}
We first verify the hypotheses needed to apply Kimura's theorem to $G$.

\begin{lemma}\label{lem:gorenstein}
Let $(R,\m,k)$ be Gorenstein, with $\m^3=0$, and suppose $R$ is not a complete intersection. Then $G=\gr_{\m}R$ is a Gorenstein local $k$-algebra, is not a complete intersection, and has embedding dimension at least three.
\end{lemma}
\begin{proof}
By \cite[Lemma 7.2(1)]{KT}, $e=\edim R\ge3$ and $\Soc R=\m^2\cong k$. Choose a generator $z$ of $\m^2$. The multiplication pairing on $\m/\m^2$ with values in $kz$ is nonsingular. Indeed, a vector in its radical lifts to an element of $\m$ annihilating $\m$, which lies in $\Soc R=\m^2$ and hence has zero degree-one class. Thus $G$ has socle $kz$, and is Gorenstein with Hilbert function $(1,e,1)$.

Write $G=k[[Y_1,\ldots,Y_e]]/I$ in its graded presentation. The space of quadratic relations has dimension $\binom{e+1}{2}-1$. As $I$ contains no linear form, these relations give linearly independent classes in $I/(Y_1,\ldots,Y_e)I$. Therefore
\[
 \mu(I)\ge\binom{e+1}{2}-1>e\qquad(e\ge3).
\]
An Artinian complete intersection in this regular local ring has a defining ideal minimally generated by $e$ elements. Hence $G$ is not a complete intersection.
\end{proof}

\begin{proof}[Proof of Theorem~\ref{thm:gorenstein}]
If $\m^3=0$, Lemma~\ref{lem:gorenstein} makes $G$ a non-complete-intersection Gorenstein local ring containing a field. Kimura's theorem \cite[Theorem 7.3(1)]{K} says that $G$ is dominant. Theorem~\ref{thm:transfer} now implies that $R$ is dominant.

Suppose $e(R)\le\codim R+2$ and the residue field is infinite. Set $d=\dim R$. Choose $x_1,\ldots,x_d\in\m$ whose initial linear forms give a homogeneous system of parameters for $\gr_{\m}R$. This is possible over an infinite field. Their ideal $Q$ is a minimal reduction of $\m$, and their classes in $\m/\m^2$ are linearly independent. Since $R$ is Cohen--Macaulay, this is a regular parameter sequence. The Artinian quotient $A=R/Q$ therefore satisfies
\[
 \ell_A(A)=e(R),\qquad \edim A=\edim R-d=\codim R.
\]
These are the usual minimal-reduction equalities, also used in \cite[proof of Theorem 7.3(2)]{K}. The ring $A$ is Gorenstein and not a complete intersection: each of these properties is equivalent across quotient by a regular sequence. If $\mathfrak a$ is its maximal ideal, then
\[
 \ell_A(\mathfrak a^2)=\ell_A(A)-1-\edim A\le1.
\]
Thus $\mathfrak a^2$ is zero or a simple $A$-module, and $\mathfrak a^3=0$. The Artinian case proves that $A$ is dominant. Applying \eqref{eq:known} successively to the parameter sequence proves dominance of $R$.

For a finite residue field, put $S=R[t]_{\m R[t]}$. The map $R\to S$ is faithfully flat and local, its maximal ideal is $\m S$, and its residue field is the infinite field $k(t)$. Flatness gives
\[
 (\m S)^j/(\m S)^{j+1}\cong(\m^j/\m^{j+1})\otimes_k k(t)
 \qquad(j\ge0).
\]
Consequently the Hilbert functions, embedding dimensions, dimensions, and multiplicities agree. The closed fiber is a field, so $S$ is Gorenstein, and $S$ is a complete intersection if and only if $R$ is; see \cite[Tags 0BJL and 09Q7]{Stacks}. The infinite-field case applies to $S$. Liu's flat-descent theorem \cite[Theorem 3.8]{L} then gives dominance of $R$.
\end{proof}

\section{Examples, verification, and scope}
The two positions of the residue characteristic in the maximal-ideal filtration both occur. Let $V=\mathbb Z_p$, let $e\ge3$, and form the algebra with basis $1,X_1,\ldots,X_e,Z$ and rules
\[
 X_iX_j=\delta_{ij}Z,\qquad X_iZ=Z^2=0.
\]
The quotients $T/(p-Z)$ and $T/(p-X_1)$ have Hilbert function $(1,e,1)$ and nonsingular multiplication pairing. They are Gorenstein and dominant. In the first, $p$ generates the socle and the characteristic is $p^2$. In the second, $p$ has nonzero degree-one class, $p^2=Z\ne0$, and the characteristic is $p^3$. Their common graded counterpart has characteristic $p$. These examples illustrate the transfer; the theorem covers arbitrary residue fields and arbitrary symmetric multiplication tensors.

The accompanying standard-library Python program constructs exact finite quotients using the lattice generated by the columns of multiplication by $p-a$. Its triangular matrix has diagonal entries $p$, giving unique representatives with every coordinate in $\{0,\ldots,p-1\}$. The program checks associativity on the integral basis, stability of the relation lattice under multiplication, the quotient cardinality, the first two maximal-ideal layers, the socle, and the position and order of $p$.

The recorded run covers 258 quotient rings, including 86 Gorenstein and 172 non-Gorenstein cases. It enumerates 16,860 elements and checks 31,698 basis associativity identities. The examples have characteristics $4,8,9,27$. These are diagnostics for the construction; the proof over an arbitrary Cohen ring is the argument of Section 2.

The construction uses $\m^3=0$ at a precise point: every triple product of positive-degree basis elements vanishes, so arbitrary lifts of the multiplication coefficients remain associative. For longer local rings, lifting a multiplication table introduces additional associativity conditions. No corresponding arbitrary-length lifting theorem is asserted. Nor does Theorem~\ref{thm:transfer} imply that every non-Gorenstein short ring is dominant.

The quantitative transfer leaves a focused further problem: determine whether the non-complete-intersection Gorenstein graded rings with Hilbert function $(1,e,1)$ are uniformly dominant, and, if so, bound their dominant indices in terms of $e$. Any such bound transfers immediately to short Gorenstein rings without a coefficient field by \eqref{eq:indices}.

\paragraph{Preparation.} This manuscript was developed with AI assistance. The proof and accompanying computations have undergone internal checks; independent expert review has not been completed.

\begin{thebibliography}{99}
\bibitem{AIS} L. L. Avramov, S. B. Iyengar, and L. M. \c Sega,
\emph{Free resolutions over short local rings}, J. Lond. Math. Soc. (2) \textbf{78} (2008), 459--476.
\url{https://arxiv.org/abs/0707.4451}.
\bibitem{K} K. Kimura, \emph{Thick subcategories over weakly symmetric algebras with radical cube zero}, preprint, 2026, version 1.
\url{https://arxiv.org/abs/2609.05781v1}.
\bibitem{KT} T. Kobayashi and R. Takahashi, \emph{On the ubiquity of uniformly dominant local rings}, preprint, 2026, version 1.
\url{https://arxiv.org/abs/2603.10810v1}.
\bibitem{L} J. Liu, \emph{On Takahashi's descent question about dominant local rings}, preprint, 2026, version 2.
\url{https://arxiv.org/abs/2608.14283v2}.
\bibitem{Stacks} The Stacks Project Authors, \emph{The Stacks Project}, Tags 032A, 09Q7, and 0BJL, accessed September 9, 2026.
\url{https://stacks.math.columbia.edu}.
\bibitem{T} R. Takahashi, \emph{Dominant local rings and subcategory classification}, Int. Math. Res. Not. IMRN \textbf{2023}, no. 9, 7259--7318.
\url{https://doi.org/10.1093/imrn/rnac053}.
\bibitem{U} R. Takahashi, \emph{Uniformly dominant local rings and Orlov spectra of singularity categories}, preprint.
\url{https://arxiv.org/abs/2412.00669}.
\end{thebibliography}
\end{document}
